\section{Zero-free consequences of the transfer principle}
\label{app:additional-potts}\label{sec:additional-potts-results}

This section collects zero-free consequences of the Potts transfer theorem.
The first subsection treats the full interval \([0,1]\); the second treats
high-temperature intervals; the third gives the girth-five consequence.
For fixed \(q,\Deg\), the full-interval results provide the zero-free input
for deterministic approximation of normalized pinned partition functions
via Barvinok interpolation and the coefficient-computation framework
of \cite{Barvinok2016,PatelRegts2017}.
The coupling inputs are recalled in \Cref{sec:soft-flip-interface} and proved in the subsequent
sections, except for the independent hard-colouring input on the critical line
\(q=11\Deg/6\).

Unless explicitly redefined, this article uses the pinning and normalized-child
conventions of \Cref{sec:potts-setup}.  For probability laws \(\nu,\nu'\) on a
finite set \(\Omega\), we use
\[
 \|\nu-\nu'\|_{\rm TV}:=\frac12\sum_{\omega\in\Omega}
       |\nu(\omega)-\nu'(\omega)|.
\]
For finite-dimensional vectors, \(\|\cdot\|_2\) is the Euclidean norm;
\(\|\cdot\|_{2\to2}\) is its induced matrix norm.  The local
\(L^2\) conventions used for the spectral-gap proof are specified in
\Cref{sec:fixed-girth-ci}.

\subsection{Zero-free neighbourhoods of \texorpdfstring{\([0,1]\)}{[0,1]}}

\begin{theorem}[Additional Potts zero-free regimes]
\label{thm:additional-potts-zf}
Let \(q,\Deg\) be integers with \(\Deg\ge2\), and assume one of the
following:
\begin{enumerate}[label=\textup{(\roman*)},leftmargin=2.2em]
  \item\label{case:intro-near-vigoda}
        \(q\ge(11/6-1/84000)\Deg\);
  \item\label{case:intro-1809}
        \(\Deg\ge125\) and \(q\ge1.809\Deg\);
  \item\label{case:intro-large-girth}
        \(\Deg\ge3\) and \(q\ge\Deg+3\).
\end{enumerate}
There exists \(\eps=\eps(q,\Deg)>0\), and in
regime~\textup{(iii)} there also exists
\(g_*=g_*(q,\Deg)\ge3\), such that the following holds.  For every finite
simple graph \(G\) of maximum degree at most \(\Deg\) and every pinning
\(\tau\), with the additional condition
\(\operatorname{girth}(G^\tau)\ge g_*\) in regime~\textup{(iii)},
\[
  \nZpin{G}{\tau}(z)\ne0
  \qquad\text{for every }z\in\mathcal U_\eps([0,1]).
\]
Consequently, \(\Zpin{G}{\tau}(z)\) has no zeros in the same neighbourhood
except possibly at \(z=0\); its order of vanishing there is
\(m_G(\tau)\).
\end{theorem}

\begin{remark}[Literature and algorithmic use of the three \texorpdfstring{$[0,1]$}{[0,1]} regimes]
Regime~\textup{(i)} continues the strict \(11\Deg/6\) coupling argument
proved in the main article \cite{ShaoShi2026ZeroFreeness}.  On the critical line, the endpoint input is the
hard-colouring coupling of Chen, Feng, Guo, Zhang, and Zou, with the
list-contraction estimates originating in the bounded-flip work of Chen,
Delcourt, Moitra, Perarnau, and Postle; see
\cite[Theorem~20 and Proposition~22]{CFFGZZ2024} and
\cite[Section~6 and Appendix~E]{CDMPP2019}.  Regime~\textup{(ii)} is the
Carlson--Vigoda \(1.809\Deg\) improvement, extended here from the hard
endpoint to the full activity interval; compare
\cite{CarlsonVigoda2024}.  Regime~\textup{(iii)} uses the large-girth
coupling mechanism of \cite{CLMM2023}, extended here to the soft activity
interval.

In contrast, the BBR result below concerns only \([x_0,1]\), not the
full activity interval; its conclusion is a high-temperature zero-free
neighbourhood.
\end{remark}

The following elementary observation isolates the finitely many integer
pairs on the critical line that occur in regime~\textup{(i)}.

\begin{lemma}[Integer reduction]
\label{lem:int-reduction}
Let \(\Deg\) and \(q\) be integers with \(3\le\Deg\le124\) and
\[
  q\ge\left(\frac{11}{6}-\frac1{84000}\right)\Deg.
\]
Then either \(q>11\Deg/6\), or
\((\Deg,q)=(6j,11j)\) for some integer \(1\le j\le20\).
\end{lemma}

\begin{proof}
The interval
\([11\Deg/6-\Deg/84000,11\Deg/6]\) has length less than \(1/6\).
The fractional part of \(11\Deg/6\) belongs to
\(\{0,1/6,\ldots,5/6\}\), so this interval contains an integer only when
\(11\Deg/6\) is itself an integer.  Since \(\gcd(11,6)=1\), this means
\(\Deg=6j\) and \(q=11j\); the bound \(\Deg\le124\) gives \(j\le20\).
\end{proof}

\begin{remark}[The hard endpoint on the critical line]
\label{rem:critical-scope}
Every integer pair with \(q=11\Deg/6\) and \(\Deg\ge2\) has
\((\Deg,q)=(6j,11j)\) for some \(j\ge1\).  Thus \(\Deg\ge6\), and
the hard-colouring input of
\cite[Theorem~20 and Proposition~22]{CFFGZZ2024} applies along the
entire critical line.  Its list-contraction input comes from
\cite[Section~6 and Appendix~E]{CDMPP2019}.
To match the pinning conventions, recall that the normalized law at
\(x=0\) is the uniform proper list-colouring law on \(G^\tau\), with
\[
  \hardlist{\tau}{u}=\{c\in\colours:\bcount_u^\tau(c)=0\},
  \qquad
  |\hardlist{\tau}{u}|-\deg_{G^\tau}(u)
  \ge q-\Deg=\frac56\Deg.
\]
The conditional weights in \cite[Section~2]{CFFGZZ2024} likewise omit
edges within the pinned set, so this correspondence allows improper
\(\tau\).  Apply its Theorem~20 to \(\tau^{r=a}\) and
\(\tau^{r=b}\), and restrict the coupling to
\(V^\tau\setminus\{r\}\), to obtain \Cref{def:potts-ci} at \(x=0\).
This supplies the independent endpoint input in
\Cref{cor:critical-line-input}; the positive-temperature estimate there
does not provide a constant uniform as \(x\downarrow0\).
\end{remark}

\begin{theorem}[Coupling inputs in the three additional Potts regimes]
\label{thm:potts-ci-regimes}
Let \(q,\Deg\) be integers with \(\Deg\ge2\), and assume one of the three
regimes in \Cref{thm:additional-potts-zf}.  In
regimes~\textup{(i)} and~\textup{(ii)}, put \(\mathcal G=\Gdeg\).  In
regime~\textup{(iii)}, let \(g_*=g_*(q,\Deg)\) be supplied by
\Cref{thm:high-girth-soft-ci} and put \(\mathcal G=\Gdegg{g_*}\).
There exists \(C_0=C_0(q,\Deg)<\infty\) such that \(\mathcal G\)
satisfies \(C_0\)-coupling independence at \(x=0\), and for every
\(\delta\in(0,1]\) there exists \(C_\delta=C_\delta(q,\Deg)<\infty\)
such that it satisfies \(C_\delta\)-coupling independence on
\([\delta,1]\).

In regimes~\textup{(ii)} and~\textup{(iii)}, one finite constant may be
used for every \(x\in[0,1]\).  The same is true in the noncritical cases
of regime~\textup{(i)}.  At the exceptional pairs
\[
  (\Deg,q)=(6j,11j),\qquad 1\le j\le20,\quad j\in\mathbb Z,
\]
the endpoint bound is supplied independently by the hard-colouring
coupling theorem, while one may take
\(C_\delta=12/(11\delta)\).  No single uniform-in-\(x\) constant is
asserted at these critical pairs.
\end{theorem}

\begin{proof}
Regime~\textup{(ii)} is \Cref{thm:cv-ci}, and regime~\textup{(iii)} is
\Cref{thm:high-girth-soft-ci}; both theorems give one constant on
\([0,1]\).

For regime~\textup{(i)}, first suppose \(\Deg\ge125\).  Since
\[
  \frac{11}{6}-\frac1{84000}-1.809=\frac{681}{28000}>0,
\]
the assertion again follows from \Cref{thm:cv-ci}.  If \(\Deg=2\), the
integer hypothesis forces \(q\ge4>11\Deg/6\), so
\Cref{thm:strict-ci} applies.  Finally, let \(3\le\Deg\le124\).
By \Cref{lem:int-reduction}, either \(q>11\Deg/6\), when
\Cref{thm:strict-ci} applies, or \((\Deg,q)=(6j,11j)\) for
\(1\le j\le20\).  In the latter case
\Cref{cor:critical-line-input} supplies both the endpoint bound and the
stated \(C_\delta\).
\end{proof}

\begin{proof}[Proof of \Cref{thm:additional-potts-zf}]
Every listed regime has \(q\ge\Deg+1\).  In regimes~\textup{(i)} and
\textup{(ii)}, apply \Cref{thm:potts-transfer} to \(\Gdeg\) using
\Cref{thm:potts-ci-regimes}.

In regime~\textup{(iii)}, the same two theorems first give a uniform radius
for every graph in \(\Gdegg{g_*}\), where \(g_*\) is the threshold supplied
by \Cref{thm:high-girth-soft-ci}.  For a pair \((G,\tau)\) with
\(\operatorname{girth}(G^\tau)\ge g_*\), apply
\Cref{lem:pinned-leaf-realization}.  The resulting graph
\(\widetilde G\) belongs to \(\Gdegg{g_*}\) and has the same normalized
pinned polynomial, so the established radius applies.  The statement for
unnormalized pinnings follows from \eqref{eq:intro-normalization}.
\end{proof}

\subsection{High-temperature zero-free neighbourhoods}

The tree contraction of Bencs, Berrekkal, and
Regts~\cite{BencsBerrekkalRegts2025Trees} gives a further region when the
number of colours may be smaller than the maximum degree.  Retaining the
unrounded contraction parameter gives the following interval.

\begin{theorem}[A BBR-certified interval at large girth]
\label{thm:intro-bbr-interval}
Let \(q,\Deg\) be integers such that
\[
  q\ge3,
  \qquad
  \frac{\Deg}{q}\ge\frac{e-1/2}{e-1}.
\]
If \((q,\Deg)=(3,4)\), put \(x_0:=3/4\).  Otherwise put
\begin{equation}
  k:=e\frac{\Deg-q/2}{\Deg-q},
  \qquad
  x_0:=1-\frac q\Deg
       \left(1-\frac{k}{\Deg}\right)^2
       \frac{\Deg-k}{\Deg-k/2}.
  \label{eq:intro-bbr-x0}
\end{equation}
Then \(0<x_0<1\), and there are
\(g_{\rm BBR}=g_{\rm BBR}(q,\Deg)\ge3\) and
\(\eps=\eps(q,\Deg)>0\) such that, for every finite simple graph \(G\)
of maximum degree at most \(\Deg\) and every pinning \(\tau\) with
\(\operatorname{girth}(G^\tau)\ge g_{\rm BBR}\),
\[
  \nZpin{G}{\tau}(z)\ne0
  \qquad\text{for every }z\in\mathcal U_\eps([x_0,1]).
\]
After decreasing \(\eps\), if necessary, so that \(\eps<x_0\), the same
zero-free statement holds for the unnormalized pinned polynomial.
\end{theorem}

\Cref{thm:intro-bbr-interval} addresses the zero-free-neighbourhood question
in \cite[Section~5]{BencsBerrekkalRegts2025Trees}.
Outside \((q,\Deg)=(3,4)\), their endpoint is obtained from
\eqref{eq:intro-bbr-x0} by replacing \(k\) with \(e^2\).
Since \(k\le e^2\) and
\(s\mapsto(1-s/\Deg)^3/(1-s/(2\Deg))\) is decreasing on
\((0,\Deg)\), the interval above contains theirs whenever the latter is
nonempty.  The exceptional pair has an empty rounded interval and is
covered separately here.
Indeed, writing \(t=\Deg-q\), the ratio hypothesis gives
\(t/q\ge1/[2(e-1)]\), and hence
\[
  k=e\left(1+\frac{q}{2t}\right)\le e^2.
\]

\begin{proof}[Proof of \Cref{thm:intro-bbr-interval}]
For \((q,\Deg)=(3,4)\), the assertion \(0<x_0<1\) is immediate.
Otherwise the ratio hypothesis and integrality imply
\(\Deg\ge q+2\): if \(\Deg=q+1\), then the ratio hypothesis forces
\(q=3\) and \(\Deg=4\).  Writing \(t=\Deg-q\ge2\), we have
\[
  \Deg(\Deg-q)-e(\Deg-q/2)
  =t(t-e)+q(t-e/2)>0,
\]
where for \(t=2\) the right side is positive already at \(q=3\), and
for \(t\ge3\) both summands are positive.  Hence \(0<k<\Deg\).  Both
factors following \(q/\Deg\) in \eqref{eq:intro-bbr-x0} lie in
\((0,1)\), so again \(0<x_0<1\).

Let \(g_{\rm BBR}\) and \(C_{\rm BBR}\) be supplied by
\Cref{thm:bbr-large-girth-ci}.  Apply
\Cref{lem:potts-positive-response} to
\(\Gdegg{g_{\rm BBR}}\), with \(\delta=x_0\) and
\(C=C_{\rm BBR}\).  This gives one radius throughout \([x_0,1]\) for
every graph in this class and every pinning.

For the free-graph formulation, apply
\Cref{lem:pinned-leaf-realization} to \((G,\tau)\).  The resulting graph
lies in \(\Gdegg{g_{\rm BBR}}\) and has the same normalized pinned
partition function, so the radius just obtained applies.  Finally shrink
it below \(x_0/2\), so the possible monomial zero of an unnormalized
pinning is outside the neighbourhood.
\end{proof}

\begin{corollary}[General-graph high-temperature intervals]
\label{cor:intro-high-temperature}
Let \(q,\Deg\) be integers with \(\Deg\ge2\), and fix \(x_*\in(0,1]\).  If
\[
  q>\frac{11}{6}(1-x_*)\Deg,
\]
then there is \(\eps=\eps(q,\Deg,x_*)>0\) such that, for every finite
simple graph \(G\) of maximum degree at most \(\Deg\) and every pinning
\(\tau\),
\[
  \nZpin{G}{\tau}(z)\ne0
  \qquad\text{for every }z\in\mathcal U_\eps([x_*,1]).
\]
After decreasing \(\eps\), if necessary, so that \(\eps<x_*\), the same
statement holds for the unnormalized pinned polynomial.
\end{corollary}

\begin{proof}
For \(x\in[x_*,1]\), \Cref{prop:full-interval-hci} gives coupling
independence with the following uniform constant, since its bound is
decreasing in \(x\):
\[
  C_*:=
  \frac{2(1-x_*)\Deg}
       {q-(11/6)(1-x_*)\Deg}.
\]
Apply \Cref{lem:potts-positive-response} to \(\Gdeg\), with
\(\delta=x_*\) and \(C=C_*\).  This gives the normalized statement.
Shrinking the radius below \(x_*\) and using
\eqref{eq:intro-normalization} gives the unnormalized statement.
\end{proof}

\subsection{Girth-five zero-free neighbourhoods at the \texorpdfstring{$(1+\delta)\Deg$}{(1+delta)Delta} threshold}
\label{app:girth5-zero-free}

Throughout this subsection, $\delta\in(0,1]$ is the slack parameter in
$q\ge(1+\delta)\Deg$.  The unrestricted girth-five coupling theorem
\Cref{thm:girth5-ci} applies to every residual instance satisfying the degree budget and makes
no local short-cycle sparsity assumption.

\begin{theorem}[Unrestricted girth-five zero-freeness]
\label{thm:unrestricted-girth5}
Fix $0<\delta\le1$, let $\Deg\ge\Deg_5(\delta)$ be as in
\Cref{thm:girth5-ci}, and let $q\ge(1+\delta)\Deg$.
There exists $\eps_5=\eps_5(q,\Deg,\delta)>0$ such that, for every
$G\in\Gdeg$ and every pinning $\tau$ with
$\operatorname{girth}(G^\tau)\ge5$,
\[
 \nZpin{G}{\tau}(z)\ne0
 \qquad\text{for every }z\in\mathcal U_{\eps_5}([0,1]).
\]
The radius is independent of $G$ and $\tau$; no bounds on local
five-cycle or six-cycle counts are assumed.
The unnormalized polynomial $\Zpin{G}{\tau}(z)$ has no zeros in this
neighbourhood except possibly at $z=0$, where its order of vanishing is
$m_G(\tau)$.  In particular, $Z_G(z)\ne0$ throughout the same
neighbourhood when $G$ itself has girth at least five.
\end{theorem}
\begin{proof}
Since $q\ge\Deg+1$, \Cref{thm:girth5-ci,thm:potts-transfer} give
one zero-free radius on the hereditary class $\Gdegg{5}$.
For the stated free-graph condition, apply
\Cref{lem:pinned-leaf-realization} to $(G,\tau)$: the resulting graph
belongs to $\Gdegg{5}$ and has the same normalized pinned polynomial.
The same radius therefore applies.  The statement for the unnormalized
polynomial follows from \eqref{eq:intro-normalization}, since the normalized
polynomial is nonzero at $z=0$.
\end{proof}

\begin{remark}[Scope and relation to the spectral-gap input]
The spectral local-to-global framework of Chen and
Liu~\cite{ChenLiu2026LocalGlobal} concerns spectral gaps and rapid mixing
on girth-five graphs.  Here the all-pinning gap is supplied by
\Cref{thm:potts-gap-girth5}; the additional coupling-independence argument
and zero-free transfer are given in \Cref{thm:girth5-ci} and above.
Thus the zero-free statement is a consequence proved here, not a
coupling-independence claim attributed to that source.
\end{remark}
